\section{Background}

\subsection{Theoretical Foundations}

The study of feature representation in neural networks builds upon theoretical work in compressed sensing \citep{donoho2006compressed} and distributed representations \citep{thorpe1989local}. Compressed sensing demonstrates how sparse signals can be recovered from lower-dimensional projections, while distributed representations show how information can be encoded across multiple units rather than localized to individual neurons.

The notion of features as fundamental units of neural computation has evolved through multiple paradigms. Early work in vision networks identified basic feature detectors \citep{olah2020zoom}, while studies of word embeddings revealed semantic directions enabling arithmetic-like operations on word vectors \citep{mikolov2013linguistic}. Three primary definitions of ``feature'' have been proposed: features as arbitrary functions of input, features as human-interpretable properties, and features as properties that sufficiently large networks dedicate neurons to representing \citep{olah2017feature}. Our work, like \citet{elhage2022toy}, adopts the third definition.

The \emph{linear representation hypothesis} posits that features correspond to directions in activation space. This hypothesis is supported by both the architectural nature of neural networks and empirical observations. Linear representations offer key advantages: they naturally emerge from standard neural computations, enable direct feature access by subsequent layers, and provide statistical efficiency through non-local generalization \citep{bengio2013representation}.

\subsection{Superposition and Polysemanticity}

\emph{Superposition} occurs when a neural network represents more features than it has dimensions by encoding features as non-orthogonal directions in activation space. This allows models to represent $n$ features in $m < n$ dimensions, at the cost of interference between features. \emph{Polysemanticity}---the phenomenon of individual neurons responding to multiple distinct features---is a consequence of superposition.

Early investigations of superposition were motivated by observations of polysemantic neurons in language models \citep{olah2020zoom}. The relationship between sparsity and superposition, noted in work on vision and natural image statistics \citep{olshausen1997sparse}, suggested that feature frequency plays a crucial role: when features are sparse (rarely active), the interference cost of superposition is low, making it advantageous to pack more features into limited dimensions.

\subsection{Related Work in Mechanistic Interpretability}

Mechanistic interpretability originated with the Distill Circuits thread \citep{cammarata2020thread}, which focused on interpreting convolutional neural networks. A framework for analyzing transformers was established by \citet{elhage2021mathematical}, providing mathematical foundations for decomposing transformer components. This enabled subsequent work such as the identification of ``induction heads'' that play a role in in-context learning \citep{olsson2022context}.

The foundational study of superposition \citep{elhage2022toy} demonstrated how toy models exhibit superposition under controlled conditions. This work was extended by \citet{henighan2023superposition}, who investigated relationships between superposition, memorization, and double descent. \citet{hubinger2023privileged} explored privileged bases in transformer residual streams, while \citet{bricken2023monosemanticity} introduced dictionary learning techniques for decomposing language models. Most recently, \citet{templeton2024scaling} demonstrated how these approaches scale to production models like Claude 3 Sonnet, showing that identified features can be used to steer model behavior---a key motivation for understanding superposition.

Architectural innovations such as Softmax Linear Units \citep{elhage2022softmax} aim to make neural networks more interpretable while maintaining performance. Our work builds on this foundation by testing whether superposition findings from toy models generalize to practical architectures.
